\documentclass[ignorenonframetext]{beamer}

\mode<article>
{
  \usepackage{fullpage}
  \usepackage{pgf}
  \usepackage{hyperref}
  \setjobnamebeamerversion{vortrag.beamer}
}

% WK OLD
% \mode<presentation>
% {
%   \usetheme{Darmstadt}
%   \useoutertheme{smoothbarssmall}
%   \usecolortheme{seahorse}
% %   \setbeamercovered{transparent}
% }
\mode<presentation> {
    \usetheme{Duesseldorf}
    \setbeamercovered{transparent}
}

\setbeamertemplate{footline}{
  {%
    \begin{beamercolorbox}[colsep=1.5pt]{upper separation line foot}
    \end{beamercolorbox}
%     \begin{beamercolorbox}[ht=2.5ex,dp=1.125ex,%
%       leftskip=.3cm,rightskip=.3cm plus1fil]{author in head/foot}%
%       \leavevmode{\usebeamerfont{author in head/foot}\insertshortauthor}%
%       \hfill%
%       {\usebeamerfont{institute in head/foot}\usebeamercolor[fg]{institute in head/foot}\insertshortinstitute}%
%     \end{beamercolorbox}%
    \begin{beamercolorbox}[ht=2.5ex,dp=1.125ex,%
      leftskip=.3cm,rightskip=.3cm plus1fil]{title in head/foot}%
      {\usebeamerfont{title in head/foot}\insertshorttitle}%
      \hfill%
      \insertframenumber% /\inserttotalframenumber
    \end{beamercolorbox}%
    \begin{beamercolorbox}[colsep=1.5pt]{lower separation line foot}
    \end{beamercolorbox}
  }
}

\usepackage[english,ngerman]{babel}
\usepackage{ngerman}
\usepackage[utf8]{inputenc}
\usepackage{times,mathptmx,eurosym}
\usepackage[arrow, matrix, curve]{xy}
\usepackage[T1]{fontenc}
\usepackage{multimedia}
\usepackage{amsfonts}
\usepackage{amssymb}
\usepackage{amstext}
\usepackage{multirow}
\usepackage{pict2e}
\usepackage{tabularx}
%\usepackage{animate}
\newcommand{\reference}[1]{{\scriptsize [#1]}}


\title[]{Eigenschaften eines Algorithmus zur Zuordnung von nicht eindeutigen Sende- und Empfangsereignissen in Ereignisprotokollen}

% \subtitle{Implicit Feedback in Wireless Multihop Networks}
%{Untertitel nur angeben, wenn es einen im Tagungsband gibt}

\author{Benito van der Zander}
% - Namen ssen in derselben Reihenfolge wie im Papier erscheinen.
% - Der \inst{?} Befehl sollte nur verwendet werden, wenn die Autoren
%   unterschiedlichen Instituten angeh�ren.

\institute%[Universit�ten Hier und Dort] % (optional, aber oft n�tig)
{Lehrstuhl für Rechnernetze\\
Heinrich-Heine-Universität D\"usseldorf}

\date{14. September 2009}

% \date[VTC Spring 2007] % (optional, sollte der abgek�rzte Konferenzname sein)
% {23 April 2007}
% \date{Universität Bonn \quad 16 August 2007}
% - Vollee oder abgek�rzter Name sind m�glich.
% - Dieser Eintrag ist nicht f�r das Publikum gedacht (das wei�
%   n�mlich, bei welcher Konferenz es ist), sondern f�r Leute, die die
%   Folien sp�ter lesen.

% \subject{Offline Time Synchronization for libpcap Logs}
% Dies wird lediglich in den PDF Informationskatalog einf�gt. Kann gut
% weggelassen werden.


% Falls eine Logodatei namens "university-logo-filename.xxx" vorhanden
% ist, wobei xxx ein von latex bzw. pdflatex lesbares Graphikformat
% ist, so kann man wie folgt ein Logo einf�gen:

% \pgfdeclareimage[width=1.5cm]{university-logo}{unilogo}
% \logo{\pgfuseimage{university-logo}}
\logo{\includegraphics[width=1.5cm]{unilogo}}
% Folgendes sollte gel�scht werden, wenn man nicht am Anfang jedes
% Unterabschnitts die Gliederung nochmal sehen m�chte.
% \AtBeginSubsection[]
% {
%   \begin{frame}<beamer>
%     \frametitle{Gliederung}
%     \tableofcontents[currentsection,currentsubsection]
%   \end{frame}
% }


% Falls Aufz�hlungen immer schrittweise gezeigt werden sollen, kann
% folgendes Kommando benutzt werden:

%\beamerdefaultoverlayspecification{<+->}

\setbeamercolor{question}{fg=black,bg=yellow}
\setbeamercolor{nextquestion}{fg=black,bg=green}




\newcommand{\vvv}[3]{\mbox{$\left( \begin{smallmatrix} #1\\ #2\\ #3 \end{smallmatrix} \right)$}}

\newcommand{\mybreak}{\par\vspace*{\baselineskip}}

\begin{document}


\begin{frame}
  \titlepage
\end{frame}


%%%%%%%%%%%%%%%%%%%%%%%%%%% EINLEITUNG
% \section*{Einleitung}

\begin{frame}
  \frametitle{Übersicht}

\tableofcontents


\end{frame}

\section{Problemstellung}
\begin{frame}
	\animate<1-10>
	Problem: Nachträgliches Berechnen der möglichen Sender für empfangene Nachrichten.\mybreak %Wie berechnet man aus aufgezeichneten Ereignissen welche Nachricht empfangen wurde?\mybreak 
	
	
			Sendende und empfangene Knoten:\mybreak
	
	\begin{center}
%	\begin{tabular}{ccccccccc}
%		A &   &  A&   &    & A  &   &   &$\ldots$\\
%		  & B &   &   &     &   & B &   & $\ldots$ \\
%		  &   &   & C &  C  &   &   &   &$\ldots$\\
%		  & D &   &   &     &   &   & D & $\ldots$\\
%	\end{tabular}
	\begin{tabular}{ccccccccccc}
		\onslide<1->\alert<1>A & \onslide<2->\alert<2>A &\onslide<3->\alert<3>b&\onslide<4->\alert<4>A&\onslide<5->\alert<5>c&     & \onslide<7->\alert<7>A  &\onslide<8->\alert<8>b  &   &\onslide<10->$\ldots$\\
		\onslide<1->\alert<1>a &               &\onslide<3->\alert<3>B&             &             &\onslide<6->\alert<6>c  &   & \onslide<8->\alert<8>B & \onslide<9->\alert<9>d & \onslide<10->$\ldots$ \\
		\onslide<1->\alert<1>a &               &\onslide<3->\alert<3>d&\onslide<4->\alert<4>a&\onslide<5->\alert<5>C&\onslide<6->\alert<6>C  &   & \onslide<8->\alert<8>b &   &\onslide<10->$\ldots$\\
		              &               &\onslide<3->\alert<3>D&\onslide<4->\alert<4>a&   &     &   &   & \onslide<9->\alert<9>D & \onslide<10->$\ldots$\\
	\end{tabular}
	\end{center}

	\onslide<1->Führen zu folgenden Logs:\mybreak
	\begin{center}
	\begin{tabular}{ccccccccccccccccccccccc}
		\onslide<1->{\alert<1>A} & \onslide<2->{\alert<2>A} & \onslide<3->{\alert<3>b} & \onslide<4->{\alert<4>A} & \onslide<5->{\alert<5>c} & \onslide<7->{\alert<7>A}  &\onslide<8->{\alert<8>b}  & \onslide<10->{$\ldots$}\\
		\onslide<1->{\alert<1>a} & \onslide<3->{\alert<3>B} & \onslide<6->{\alert<6>c} & \onslide<8->{\alert<8>B} & \onslide<9->{\alert<9>d} &  \onslide<10->{$\ldots$} \\
		\onslide<1->{\alert<1>a} & \onslide<3->{\alert<3>d} & \onslide<4->{\alert<4>a} & \onslide<5->{\alert<5>C} & \onslide<6->{\alert<6>C} & \onslide<8->{\alert<8>b} &  \onslide<10->{$\ldots$}\\
		\onslide<3->{\alert<3>D} & \onslide<4->{\alert<4>a} & \onslide<9->{\alert<9>D} & \onslide<10->{$\ldots$ }\\
	\end{tabular}
	\end{center}
				
	
\end{frame}
\begin{frame}
	Mögliche Fragen zu diesen Logs:\mybreak
	
	\begin{description}
		%Für ein Ereignis $r \in R$, was ist die Menge der möglichen Sendeereignisse, also für welche $s$ existiert ein  konsistenter GA $A$ mit $A(r) = s$
		
		\item[SRC: ] Kann eine konsistente Zuordnung gefunden werden?%Existiert eine konsistente GA
		
		\item[uSRC:] Kann eine konsistente und eindeutige Zuordnung gefunden werden?
		
		\item[pSRC:] Was ist für ein gegebene Empfangsereignis die Menge der möglichen Sendeereignisse?
		
	\end{description}
	
\end{frame}
\section{Komplexität}
\subsection{3-Exact-Cover (X3C)}
\begin{frame}
	\animate<1-3>
	Existiert eine exakte Mengenüberdeckung?\mybreak
	
	Es ist eine Menge gegeben: \[M = \{A, B, C, D, E, F\}\]
	
	\pause

	Zusammen mit dreielementigen Untermengen:
	\begin{center}
		\begin{tabular}{cccccccc}
	$M_1 = \{$& $A$ & $B$ & $C$ &   & & &$\}$\\
 	\color<3>{red}$M_2 = \{$& \color<3>{red}$A$ &   & \color<3>{red}$C$ & \color<3>{red}$D$ && &\color<3>{red}$\}$\\
	$M_3 = \{$&   & $B$ & $C$ & $D$ && &$\}$\\
	 $M_4 = \{$&   &   & $C$ & $D$ & $E$ && $\}$\\
	\color<3>{red}$M_5 = \{$&   & \color<3>{red}$ B $  &   &   &\color<3>{red}$ E$ & \color<3>{red}$F$& \color<3>{red}$\}$
		\end{tabular}
	\end{center}


\end{frame}

\subsection{Unique Send-Receive-Correlation (uSRC)}

\begin{frame}
	\animate<1-6>
Umwandeln des X3C-Problems in ein uSRC-Problems:\mybreak

Sendelog:
\only<1-2>{\[L_0 = \{A, B, C, D, E, F\}\]}
\only<3-4>{\[L_0 = \{A, B, C, D, E, F, \quad A, B, B, C,C,C,D,D,E\}\]}
\only<5-6>{\[L_0 = \{A, B, C, D, E, F, X, X, X, X, X, X, A, B, B, C,C,C,D,D,E\}\]}

Empfangslogs:
	\begin{center}
		\begin{tabular}{ccccccccc}
			$L_1 = \{$& $a$ & $b$ & $c$ &   & & &\only<5-6>{$x$} & $\}$\\
			$L_2 = \{$& $a$ &   & $c$ & $d$ && &\only<5-6>{$x$}&$\}$\\
			$L_3 = \{$&   & $b$ & $c$ & $d$ && &\only<5-6>{$x$}&$\}$\\
			$L_4 = \{$&   &   & $c$ & $d$ & $e$ &&\only<5-6>{$x$}& $\}$\\
			$L_5 = \{$&   & $ b $  &   &   &$ e$ & $f$&\only<5-6>{$x$}& $\}$
		\end{tabular}
	\end{center}

\only<2>{\color{red}Problem: Zu wenig Sender!}
\only<4>{\color{red}Problem: Überlappungen!}
\only<6>{\color{orange}uSRC ist NP-vollständig!}
\end{frame}

\subsection{Send-Receive-Correlation (SRC)}
\begin{frame}
	\animate<1-4>
	Was ist mit SRC?\mybreak

\pause

	Betrachte:
	
	\begin{center}
\begin{tabular}{cccc}
	A&b&$\square$&A\\
	B&a&$\boxdot$&B\\
\end{tabular}
	\end{center}

\only<3>{
Zwei mögliche Zuordnungen:\mybreak
\begin{tabular}{ccccc}
	&$A$&$b$        &$\square$&$A$ \\
	&$\updownarrow$\\
	$B$&$a$&$\boxdot$&$B$        &  \\
\end{tabular}
\hspace{1cm}\vline\hspace{1cm}
\begin{tabular}{cccccc}
	$A$&$b$&$\square$&$A$\\
	& &         &$\updownarrow$\\
	& &  $B$      &$a$&$\boxdot$&$B$\\
\end{tabular}\mybreak

}

\only<4>{
Zwei mögliche Zuordnungen:\mybreak
\begin{tabular}{cccccc}
	&$A$ &         &$b$        &$\square$&$A$ \\
	&  $\updownarrow$ & &$\updownarrow $\\
	$B$&$a$&$\boxdot$&$B$        &  \\
\end{tabular}
\hspace{1cm}\vline\hspace{1cm}
\begin{tabular}{cccccc}
	$A$& $b$&$\square$&$A$\\
	&   $\updownarrow $      & &$\updownarrow$\\
	&  $B$ &     &$a$&$\boxdot$&$B$\\
\end{tabular}\mybreak

$\boxdot$ und $\square$ sind nicht gleichzeitig!
}


\end{frame}

\begin{frame}
	Damit reduzieren wir uSRC auf SRC:\mybreak

\only<1>{	
Sendelog:
\[L_0 = \{A, B, C, D, E, F, X, X, X, X, X, X, A, B, B, C,C,C,D,D,E\}\]}

Empfangslogs:
\begin{center}
	\only<1>{
	\begin{tabular}{ccccccccc}
		$L_1 = \{$& $a$ & $b$ & $c$ &   & & &{$x$} & $  $$\}$\\
		$L_2 = \{$& $a$ &   & $c$ & $d$ && &{$x$}&$\}$\\
		$L_3 = \{$&   & $b$ & $c$ & $d$ && &{$x$}&$\}$\\
		$L_4 = \{$&   &   & $c$ & $d$ & $e$ &&{$x$}& $\}$\\
		$L_5 = \{$&   & $ b $  &   &   &$ e$ & $f$&{$x$}& $\}$
	\end{tabular}}
%	\only<1>{\begin{tabular}{ccccccccc}
%		$L_1 = \{$& $a$ & $b$ & $c$ &   & & &{$x$} & $  $$\}$\\
%		$L_2 = \{$& $a$ &   & $c$ & $d$ && &{$x$}&$\}$\\
%		$L_3 = \{$&   & $b$ & $c$ & $d$ && &{$x$}&$\}$\\
%		$L_4 = \{$&   &   & $c$ & $d$ & $e$ &&{$x$}& $\}$\\
%		$L_5 = \{$&   & $ b $  &   &   &$ e$ & $f$&{$x$}& $\}$
%	\end{tabular}}
		\only<2-3>{
		\begin{tabular}{cccccccccccc}
		$L_1 = \{$&$A_{12}$ & $a_{21}$ & $a$ & $A_{12}$ & $b$ & $c$ &   & & &{$x$} & $\}$\\
		$L_2 = \{$&$A_{21}$ & $a_{12}$ & $a$ &$A_{21}$ &  & $c$ &  $d$ & & &{$x$} & $\}$\\
		$L_3 = \{$&&&   && $b$ & $c$ & $d$ && &{$x$}&$\}$\\
		$L_4 = \{$&&&   &&   & $c$ & $d$ & $e$ &&{$x$}& $\}$\\
		$L_5 = \{$&&&   &&$ b $  &   &   &$ e$ & $f$&{$x$}& $\}$
	\end{tabular}}
	\only<4-5>{
		\begin{tabular}{ccccccccccccccc}
			$L_1 = \{$&$A_{12}$&$A_{12}$ & $a_{21}$ & $a$ & $A_{12}$ & $b$ & $c$ &   & & &{$x$} &$A_{12}$& $\}$\\
			$L_2 = \{$&$A_{22}$&$A_{21}$ & $a_{12}$ & $a$ &$A_{21}$ &  & $c$ & $d$  & & &{$x$} &$A_{21}$& $\}$\\
			$L_3 = \{$&&&   &&& $b$ & $c$ & $d$ && &{$x$}&&$\}$\\
			$L_4 = \{$&&&   &&&   & $c$ & $d$ & $e$ &&{$x$}&& $\}$\\
			$L_5 = \{$&&&   &&&$ b $  &   &   &$ e$ & $f$&{$x$}&& $\}$
		\end{tabular}}
	\only<6-7>{
		\begin{tabular}{ccccccccccccccc}
			$\ldots B_{13}B_{15}$ & $\ldots B_{13}$ & $B_{15}$ & $b_{31}$ & $b_{51}$ & $b$ & $B_{13}$ & $B_{15} \ldots$ & $B_{13}B_{15}\ldots$\\
			&&&&$\ldots$\\
			$\ldots B_{31}B_{35}$ & $\ldots B_{31}$ & $B_{35}$ & $b_{13}$ & $b_{53}$ & $b$ & $B_{31}$ & $B_{35} \ldots$ & $B_{31}B_{35}\ldots$\\
			&&&&$\ldots$\\
			$\ldots B_{51}B_{53}$ & $\ldots B_{51}$ & $B_{53}$ & $b_{15}$ & $b_{35}$ & $b$ & $B_{51}$ & $B_{53} \ldots$ & $B_{51}B_{53}\ldots$\\
			
		\end{tabular}}
	\end{center}

	\only<3>{\color{red}Problem: Können die Sendeereignisse  zugewiesen werden?}
	\only<5>{\color{red}Problem: Mehr als zwei Ereignisse}
	\only<7>{\color{orange}Somit ist SRC NP-vollständig!}
		
\end{frame}

%\begin{frame}
%	\newcount\opaqueness
%	\animate<2-10>
%	
%	\animatevalue<1-10>{\opaqueness}{100}{0}
%	abc
%	\begin{colormixin}{\the\opaqueness!averagebackgroundcolor}
%		\frametitle{Fadeout Frame}
%		This text (and all other frame content) will fade out when the
%		second slide is shown. This even works with
%		{\color{green!90!black}colored} \alert{text}.
%	\end{colormixin}
%\end{frame}

\section{Intervallbasierte Zuordnung und Heuristik}}
\subsection{Intervalle}
\begin{frame}
	
	\animate<1-2>
	
	Bestehen die Lösungen aus Intervallen?\mybreak
	
	Betrachte:\mybreak
	
	\begin{tabular}{cccccc}
		B & c & a & B &   & \\
		C & A & C & b & C & A
	\end{tabular} \mybreak
	
	\pause
	
	Mögliche Lösungen:\mybreak
	
	
	
	\begin{columns}
		\begin{column}{5cm}
			\begin{tabular}{cccccccc}
				B & c & a & & B \\
				& $\uparrow$ & $\uparrow$ && $\downarrow$ \\
				& C & A & C & b & C & A
			\end{tabular} 
		\end{column}


		\begin{column}{5cm}
			\begin{tabular}{cccccccc}
				&   &   & B & c & a & B \\
				&   &   & $\downarrow$ & $\uparrow$ & $\uparrow$ & \\
				C & A & C & b & C & A
			\end{tabular} 
		\end{column}
	\end{columns}
	
	\mybreak Das Sendeintervall für c enthält sämtliche gesendeten C.
	
%	\onslide<3>
%	\begin{figure}[h]
%		\begin{minipage}[h]{6cm}
%			\begin{tabular}{ccccccc}
%				& B & c &  a  & B \\
%				&   & $\uparrow$ &  ???? & $\downarrow$\\
%				C & A & C &     & b & C & A
%			\end{tabular} 
%		\end{minipage}
%		\hspace{1cm}\vline\hspace{1cm}
%		\begin{minipage}[h]{6cm}
%			\begin{tabular}{cccccccc}
%				& B & c &   &   & a  & B \\
%				&   & $\uparrow$ &  ???? && $\uparrow$ & \\
%				C & A & C & b & C & A
%			\end{tabular} 
%		\end{minipage}
%	\end{figure}
%	
\end{frame}

\begin{frame}
	Aber ist das mittlere C ein gültiger Sender?\mybreak\mybreak


	\begin{columns}
		\begin{column}{6cm}
			\begin{tabular}{ccccccc}
				& B & c &  a  & B \\
				&   & $\uparrow$ &  ???? & $\downarrow$\\
				C & A & C &     & b & C & A
			\end{tabular} 
		\end{column}

		\begin{column}{6cm}
			\begin{tabular}{cccccccc}
				& B & c &   &   & a  & B \\
				&   & $\uparrow$ &  ???? && $\uparrow$ & \\
				C & A & C & b & C & A
			\end{tabular} 
		\end{column}
	\end{columns}\mybreak\mybreak

	Offensichtlich nicht!
\end{frame}

\subsection{Heuristik}
\begin{frame}
	Findet die Heuristik trotzdem gute Intervalle?\mybreak
	
	Betrachte:\mybreak
	
	\begin{center}
		\begin{tabular}{cccccc}
		A & b & c & A\\
		B & a & B & c & B\\
		C & a & b & C
	\end{tabular}
\end{center} \mybreak
	
	Diese Eingabe ist lösbar:\mybreak
	
	\center{\begin{tabular}{cccccc}
		& A & b & c & A\\
		& $\updownarrow$ & $\updownarrow$ & $\updownarrow$\\
		B & a & B & c & B\\
		& $\updownarrow$ & $\updownarrow$ & $\updownarrow$\\
		C & a & b & C
	\end{tabular} }
	
\end{frame}

\begin{frame}
	Die Ausgabe lautet:\mybreak
	
	\setlength{\unitlength}{0.7pt}
	\begin{picture}(380,206)(30,196)
		\definecolor{R107G107B153}{rgb}{0.41960785,0.41960785,0.6}
		\definecolor{R153G107B107}{rgb}{0.6,0.41960785,0.41960785}
		\definecolor{R107G153B107}{rgb}{0.41960785,0.6,0.41960785}
		\definecolor{R0G128B0}{rgb}{0.0,0.5019608,0.0}
		\definecolor{R0G179B0}{rgb}{0.0,0.7019608,0.0}
		\definecolor{R0G0B128}{rgb}{0.0,0.0,0.5019608}
		\definecolor{R0G0B179}{rgb}{0.0,0.0,0.7019608}
		\definecolor{R128G0B0}{rgb}{0.5019608,0.0,0.0}
		\definecolor{R179G0B0}{rgb}{0.7019608,0.0,0.0}
		\put(60,322){\color{R107G107B153}\linethickness{1.0pt}\vector(80,30){80}}
		\put(140,352){\color{R107G107B153}\linethickness{1.0pt}\vector(240,-30){240}}
		\put(60,244){\color{R153G107B107}\linethickness{1.0pt}\vector(160,108){160}}
		\put(220,352){\color{R153G107B107}\linethickness{1.0pt}\vector(80,-108){80}}
		\put(60,352){\color{R107G153B107}\linethickness{1.0pt}\vector(80,-30){80}}
		\put(140,322){\color{R107G153B107}\linethickness{1.0pt}\vector(160,30){160}}
		\put(60,244){\color{R153G107B107}\linethickness{1.0pt}\vector(240,30){240}}
		\put(300,274){\color{R153G107B107}\linethickness{1.0pt}\vector(0,-30){30}}
		\put(60,352){\color{R107G153B107}\linethickness{1.0pt}\vector(80,-108){80}}
		\put(140,244){\color{R107G153B107}\linethickness{1.0pt}\vector(160,108){160}}
		\put(60,274){\color{R107G107B153}\linethickness{1.0pt}\vector(160,-30){160}}
		\put(220,244){\color{R107G107B153}\linethickness{1.0pt}\vector(160,30){160}}
		\put(50,388){\makebox(20,20)[b]{\fontseries{b}\fontsize{12.0}{14}\selectfont  A}}
		\put(37,379){\makebox(20,20)[br]{\fontsize{10.0}{12}\selectfont  0}}
		\put(37,367){\makebox(20,20)[br]{\fontsize{10.0}{12}\selectfont  0}}
		\put(37,355){\makebox(20,20)[br]{\fontsize{10.0}{12}\selectfont  0}}
		\put(67,379){\makebox(20,20)[br]{\fontsize{10.0}{12}\selectfont  0}}
		\put(67,367){\makebox(20,20)[br]{\fontsize{10.0}{12}\selectfont  1}}
		\put(67,355){\makebox(20,20)[br]{\fontsize{10.0}{12}\selectfont  1}}
		\put(60,385){\color{R0G179B0}\linethickness{2.0pt}\line(0,-30){30}}
		\put(30,388){\color{R0G179B0}\linethickness{2.0pt}\line(0,-36){36}}
		\put(30,388){\color{R0G179B0}\linethickness{2.0pt}\line(5,0){5}}
		\put(30,352){\color{R0G179B0}\linethickness{2.0pt}\line(5,0){5}}
		\put(90,388){\color{R0G179B0}\linethickness{2.0pt}\line(0,-36){36}}
		\put(90,388){\color{R0G179B0}\linethickness{2.0pt}\line(-5,0){5}}
		\put(90,352){\color{R0G179B0}\linethickness{2.0pt}\line(-5,0){5}}
		\put(130,388){\makebox(20,20)[b]{\fontseries{b}\fontsize{12.0}{14}\selectfont  b}}
		\put(117,379){\makebox(20,20)[br]{\fontsize{10.0}{12}\selectfont  1}}
		\put(117,367){\makebox(20,20)[br]{\fontsize{10.0}{12}\selectfont  0}}
		\put(117,355){\makebox(20,20)[br]{\fontsize{10.0}{12}\selectfont  0}}
		\put(147,379){\makebox(20,20)[br]{\fontsize{10.0}{12}\selectfont  1}}
		\put(147,367){\makebox(20,20)[br]{\fontsize{10.0}{12}\selectfont  4}}
		\put(147,355){\makebox(20,20)[br]{\fontsize{10.0}{12}\selectfont  3}}
		\put(140,385){\color{R0G0B179}\linethickness{2.0pt}\line(0,-30){30}}
		\put(110,388){\color{R0G0B179}\linethickness{2.0pt}\line(0,-36){36}}
		\put(110,388){\color{R0G0B179}\linethickness{2.0pt}\line(5,0){5}}
		\put(110,352){\color{R0G0B179}\linethickness{2.0pt}\line(5,0){5}}
		\put(170,388){\color{R0G0B179}\linethickness{2.0pt}\line(0,-36){36}}
		\put(170,388){\color{R0G0B179}\linethickness{2.0pt}\line(-5,0){5}}
		\put(170,352){\color{R0G0B179}\linethickness{2.0pt}\line(-5,0){5}}
		\put(210,388){\makebox(20,20)[b]{\fontseries{b}\fontsize{12.0}{14}\selectfont  c}}
		\put(197,379){\makebox(20,20)[br]{\fontsize{10.0}{12}\selectfont  2}}
		\put(197,367){\makebox(20,20)[br]{\fontsize{10.0}{12}\selectfont  0}}
		\put(197,355){\makebox(20,20)[br]{\fontsize{10.0}{12}\selectfont  0}}
		\put(227,379){\makebox(20,20)[br]{\fontsize{10.0}{12}\selectfont  2}}
		\put(227,367){\makebox(20,20)[br]{\fontsize{10.0}{12}\selectfont  4}}
		\put(227,355){\makebox(20,20)[br]{\fontsize{10.0}{12}\selectfont  3}}
		\put(220,385){\color{R179G0B0}\linethickness{2.0pt}\line(0,-30){30}}
		\put(190,388){\color{R179G0B0}\linethickness{2.0pt}\line(0,-36){36}}
		\put(190,388){\color{R179G0B0}\linethickness{2.0pt}\line(5,0){5}}
		\put(190,352){\color{R179G0B0}\linethickness{2.0pt}\line(5,0){5}}
		\put(250,388){\color{R179G0B0}\linethickness{2.0pt}\line(0,-36){36}}
		\put(250,388){\color{R179G0B0}\linethickness{2.0pt}\line(-5,0){5}}
		\put(250,352){\color{R179G0B0}\linethickness{2.0pt}\line(-5,0){5}}
		\put(290,388){\makebox(20,20)[b]{\fontseries{b}\fontsize{12.0}{14}\selectfont  A}}
		\put(277,379){\makebox(20,20)[br]{\fontsize{10.0}{12}\selectfont  3}}
		\put(277,367){\makebox(20,20)[br]{\fontsize{10.0}{12}\selectfont  1}}
		\put(277,355){\makebox(20,20)[br]{\fontsize{10.0}{12}\selectfont  1}}
		\put(307,379){\makebox(20,20)[br]{\fontsize{10.0}{12}\selectfont  3}}
		\put(307,367){\makebox(20,20)[br]{\fontsize{10.0}{12}\selectfont  4}}
		\put(307,355){\makebox(20,20)[br]{\fontsize{10.0}{12}\selectfont  3}}
		\put(300,385){\color{R0G179B0}\linethickness{2.0pt}\line(0,-30){30}}
		\put(270,388){\color{R0G179B0}\linethickness{2.0pt}\line(0,-36){36}}
		\put(270,388){\color{R0G179B0}\linethickness{2.0pt}\line(5,0){5}}
		\put(270,352){\color{R0G179B0}\linethickness{2.0pt}\line(5,0){5}}
		\put(330,388){\color{R0G179B0}\linethickness{2.0pt}\line(0,-36){36}}
		\put(330,388){\color{R0G179B0}\linethickness{2.0pt}\line(-5,0){5}}
		\put(330,352){\color{R0G179B0}\linethickness{2.0pt}\line(-5,0){5}}
		\put(50,310){\makebox(20,20)[b]{\fontseries{b}\fontsize{12.0}{14}\selectfont  B}}
		\put(37,301){\makebox(20,20)[br]{\fontsize{10.0}{12}\selectfont  0}}
		\put(37,289){\makebox(20,20)[br]{\fontsize{10.0}{12}\selectfont  0}}
		\put(37,277){\makebox(20,20)[br]{\fontsize{10.0}{12}\selectfont  0}}
		\put(67,301){\makebox(20,20)[br]{\fontsize{10.0}{12}\selectfont  1}}
		\put(67,289){\makebox(20,20)[br]{\fontsize{10.0}{12}\selectfont  0}}
		\put(67,277){\makebox(20,20)[br]{\fontsize{10.0}{12}\selectfont  2}}
		\put(60,307){\color{R0G0B179}\linethickness{2.0pt}\line(0,-30){30}}
		\put(30,310){\color{R0G0B179}\linethickness{2.0pt}\line(0,-36){36}}
		\put(30,310){\color{R0G0B179}\linethickness{2.0pt}\line(5,0){5}}
		\put(30,274){\color{R0G0B179}\linethickness{2.0pt}\line(5,0){5}}
		\put(90,310){\color{R0G0B179}\linethickness{2.0pt}\line(0,-36){36}}
		\put(90,310){\color{R0G0B179}\linethickness{2.0pt}\line(-5,0){5}}
		\put(90,274){\color{R0G0B179}\linethickness{2.0pt}\line(-5,0){5}}
		\put(130,310){\makebox(20,20)[b]{\fontseries{b}\fontsize{12.0}{14}\selectfont  a}}
		\put(117,301){\makebox(20,20)[br]{\fontsize{10.0}{12}\selectfont  0}}
		\put(117,289){\makebox(20,20)[br]{\fontsize{10.0}{12}\selectfont  1}}
		\put(117,277){\makebox(20,20)[br]{\fontsize{10.0}{12}\selectfont  0}}
		\put(147,301){\makebox(20,20)[br]{\fontsize{10.0}{12}\selectfont  3}}
		\put(147,289){\makebox(20,20)[br]{\fontsize{10.0}{12}\selectfont  1}}
		\put(147,277){\makebox(20,20)[br]{\fontsize{10.0}{12}\selectfont  3}}
		\put(140,307){\color{R0G179B0}\linethickness{2.0pt}\line(0,-30){30}}
		\put(110,310){\color{R0G179B0}\linethickness{2.0pt}\line(0,-36){36}}
		\put(110,310){\color{R0G179B0}\linethickness{2.0pt}\line(5,0){5}}
		\put(110,274){\color{R0G179B0}\linethickness{2.0pt}\line(5,0){5}}
		\put(170,310){\color{R0G179B0}\linethickness{2.0pt}\line(0,-36){36}}
		\put(170,310){\color{R0G179B0}\linethickness{2.0pt}\line(-5,0){5}}
		\put(170,274){\color{R0G179B0}\linethickness{2.0pt}\line(-5,0){5}}
		\put(210,310){\makebox(20,20)[b]{\fontseries{b}\fontsize{12.0}{14}\selectfont  B}}
		\put(197,301){\makebox(20,20)[br]{\fontsize{10.0}{12}\selectfont  0}}
		\put(197,289){\makebox(20,20)[br]{\fontsize{10.0}{12}\selectfont  2}}
		\put(197,277){\makebox(20,20)[br]{\fontsize{10.0}{12}\selectfont  0}}
		\put(227,301){\makebox(20,20)[br]{\fontsize{10.0}{12}\selectfont  3}}
		\put(227,289){\makebox(20,20)[br]{\fontsize{10.0}{12}\selectfont  2}}
		\put(227,277){\makebox(20,20)[br]{\fontsize{10.0}{12}\selectfont  3}}
		\put(220,307){\color{R0G0B179}\linethickness{2.0pt}\line(0,-30){30}}
		\put(190,310){\color{R0G0B179}\linethickness{2.0pt}\line(0,-36){36}}
		\put(190,310){\color{R0G0B179}\linethickness{2.0pt}\line(5,0){5}}
		\put(190,274){\color{R0G0B179}\linethickness{2.0pt}\line(5,0){5}}
		\put(250,310){\color{R0G0B179}\linethickness{2.0pt}\line(0,-36){36}}
		\put(250,310){\color{R0G0B179}\linethickness{2.0pt}\line(-5,0){5}}
		\put(250,274){\color{R0G0B179}\linethickness{2.0pt}\line(-5,0){5}}
		\put(290,310){\makebox(20,20)[b]{\fontseries{b}\fontsize{12.0}{14}\selectfont  c}}
		\put(277,301){\makebox(20,20)[br]{\fontsize{10.0}{12}\selectfont  0}}
		\put(277,289){\makebox(20,20)[br]{\fontsize{10.0}{12}\selectfont  3}}
		\put(277,277){\makebox(20,20)[br]{\fontsize{10.0}{12}\selectfont  0}}
		\put(307,301){\makebox(20,20)[br]{\fontsize{10.0}{12}\selectfont  3}}
		\put(307,289){\makebox(20,20)[br]{\fontsize{10.0}{12}\selectfont  3}}
		\put(307,277){\makebox(20,20)[br]{\fontsize{10.0}{12}\selectfont  3}}
		\put(300,307){\color{R179G0B0}\linethickness{2.0pt}\line(0,-30){30}}
		\put(270,310){\color{R179G0B0}\linethickness{2.0pt}\line(0,-36){36}}
		\put(270,310){\color{R179G0B0}\linethickness{2.0pt}\line(5,0){5}}
		\put(270,274){\color{R179G0B0}\linethickness{2.0pt}\line(5,0){5}}
		\put(330,310){\color{R179G0B0}\linethickness{2.0pt}\line(0,-36){36}}
		\put(330,310){\color{R179G0B0}\linethickness{2.0pt}\line(-5,0){5}}
		\put(330,274){\color{R179G0B0}\linethickness{2.0pt}\line(-5,0){5}}
		\put(370,310){\makebox(20,20)[b]{\fontseries{b}\fontsize{12.0}{14}\selectfont  B}}
		\put(357,301){\makebox(20,20)[br]{\fontsize{10.0}{12}\selectfont  1}}
		\put(357,289){\makebox(20,20)[br]{\fontsize{10.0}{12}\selectfont  4}}
		\put(357,277){\makebox(20,20)[br]{\fontsize{10.0}{12}\selectfont  2}}
		\put(387,301){\makebox(20,20)[br]{\fontsize{10.0}{12}\selectfont  3}}
		\put(387,289){\makebox(20,20)[br]{\fontsize{10.0}{12}\selectfont  4}}
		\put(387,277){\makebox(20,20)[br]{\fontsize{10.0}{12}\selectfont  3}}
		\put(380,307){\color{R0G0B179}\linethickness{2.0pt}\line(0,-30){30}}
		\put(350,310){\color{R0G0B179}\linethickness{2.0pt}\line(0,-36){36}}
		\put(350,310){\color{R0G0B179}\linethickness{2.0pt}\line(5,0){5}}
		\put(350,274){\color{R0G0B179}\linethickness{2.0pt}\line(5,0){5}}
		\put(410,310){\color{R0G0B179}\linethickness{2.0pt}\line(0,-36){36}}
		\put(410,310){\color{R0G0B179}\linethickness{2.0pt}\line(-5,0){5}}
		\put(410,274){\color{R0G0B179}\linethickness{2.0pt}\line(-5,0){5}}
		\put(50,232){\makebox(20,20)[b]{\fontseries{b}\fontsize{12.0}{14}\selectfont  C}}
		\put(37,223){\makebox(20,20)[br]{\fontsize{10.0}{12}\selectfont  0}}
		\put(37,211){\makebox(20,20)[br]{\fontsize{10.0}{12}\selectfont  0}}
		\put(37,199){\makebox(20,20)[br]{\fontsize{10.0}{12}\selectfont  0}}
		\put(67,223){\makebox(20,20)[br]{\fontsize{10.0}{12}\selectfont  2}}
		\put(67,211){\makebox(20,20)[br]{\fontsize{10.0}{12}\selectfont  3}}
		\put(67,199){\makebox(20,20)[br]{\fontsize{10.0}{12}\selectfont  0}}
		\put(60,229){\color{R179G0B0}\linethickness{2.0pt}\line(0,-30){30}}
		\put(30,232){\color{R179G0B0}\linethickness{2.0pt}\line(0,-36){36}}
		\put(30,232){\color{R179G0B0}\linethickness{2.0pt}\line(5,0){5}}
		\put(30,196){\color{R179G0B0}\linethickness{2.0pt}\line(5,0){5}}
		\put(90,232){\color{R179G0B0}\linethickness{2.0pt}\line(0,-36){36}}
		\put(90,232){\color{R179G0B0}\linethickness{2.0pt}\line(-5,0){5}}
		\put(90,196){\color{R179G0B0}\linethickness{2.0pt}\line(-5,0){5}}
		\put(130,232){\makebox(20,20)[b]{\fontseries{b}\fontsize{12.0}{14}\selectfont  a}}
		\put(117,223){\makebox(20,20)[br]{\fontsize{10.0}{12}\selectfont  0}}
		\put(117,211){\makebox(20,20)[br]{\fontsize{10.0}{12}\selectfont  0}}
		\put(117,199){\makebox(20,20)[br]{\fontsize{10.0}{12}\selectfont  1}}
		\put(147,223){\makebox(20,20)[br]{\fontsize{10.0}{12}\selectfont  3}}
		\put(147,211){\makebox(20,20)[br]{\fontsize{10.0}{12}\selectfont  4}}
		\put(147,199){\makebox(20,20)[br]{\fontsize{10.0}{12}\selectfont  1}}
		\put(140,229){\color{R0G179B0}\linethickness{2.0pt}\line(0,-30){30}}
		\put(110,232){\color{R0G179B0}\linethickness{2.0pt}\line(0,-36){36}}
		\put(110,232){\color{R0G179B0}\linethickness{2.0pt}\line(5,0){5}}
		\put(110,196){\color{R0G179B0}\linethickness{2.0pt}\line(5,0){5}}
		\put(170,232){\color{R0G179B0}\linethickness{2.0pt}\line(0,-36){36}}
		\put(170,232){\color{R0G179B0}\linethickness{2.0pt}\line(-5,0){5}}
		\put(170,196){\color{R0G179B0}\linethickness{2.0pt}\line(-5,0){5}}
		\put(210,232){\makebox(20,20)[b]{\fontseries{b}\fontsize{12.0}{14}\selectfont  b}}
		\put(197,223){\makebox(20,20)[br]{\fontsize{10.0}{12}\selectfont  0}}
		\put(197,211){\makebox(20,20)[br]{\fontsize{10.0}{12}\selectfont  0}}
		\put(197,199){\makebox(20,20)[br]{\fontsize{10.0}{12}\selectfont  2}}
		\put(227,223){\makebox(20,20)[br]{\fontsize{10.0}{12}\selectfont  3}}
		\put(227,211){\makebox(20,20)[br]{\fontsize{10.0}{12}\selectfont  4}}
		\put(227,199){\makebox(20,20)[br]{\fontsize{10.0}{12}\selectfont  2}}
		\put(220,229){\color{R0G0B179}\linethickness{2.0pt}\line(0,-30){30}}
		\put(190,232){\color{R0G0B179}\linethickness{2.0pt}\line(0,-36){36}}
		\put(190,232){\color{R0G0B179}\linethickness{2.0pt}\line(5,0){5}}
		\put(190,196){\color{R0G0B179}\linethickness{2.0pt}\line(5,0){5}}
		\put(250,232){\color{R0G0B179}\linethickness{2.0pt}\line(0,-36){36}}
		\put(250,232){\color{R0G0B179}\linethickness{2.0pt}\line(-5,0){5}}
		\put(250,196){\color{R0G0B179}\linethickness{2.0pt}\line(-5,0){5}}
		\put(290,232){\makebox(20,20)[b]{\fontseries{b}\fontsize{12.0}{14}\selectfont  C}}
		\put(277,223){\makebox(20,20)[br]{\fontsize{10.0}{12}\selectfont  2}}
		\put(277,211){\makebox(20,20)[br]{\fontsize{10.0}{12}\selectfont  3}}
		\put(277,199){\makebox(20,20)[br]{\fontsize{10.0}{12}\selectfont  3}}
		\put(307,223){\makebox(20,20)[br]{\fontsize{10.0}{12}\selectfont  3}}
		\put(307,211){\makebox(20,20)[br]{\fontsize{10.0}{12}\selectfont  4}}
		\put(307,199){\makebox(20,20)[br]{\fontsize{10.0}{12}\selectfont  3}}
		\put(300,229){\color{R179G0B0}\linethickness{2.0pt}\line(0,-30){30}}
		\put(270,232){\color{R179G0B0}\linethickness{2.0pt}\line(0,-36){36}}
		\put(270,232){\color{R179G0B0}\linethickness{2.0pt}\line(5,0){5}}
		\put(270,196){\color{R179G0B0}\linethickness{2.0pt}\line(5,0){5}}
		\put(330,232){\color{R179G0B0}\linethickness{2.0pt}\line(0,-36){36}}
		\put(330,232){\color{R179G0B0}\linethickness{2.0pt}\line(-5,0){5}}
		\put(330,196){\color{R179G0B0}\linethickness{2.0pt}\line(-5,0){5}}
	\end{picture}
	
\end{frame}

\begin{frame}
	\animate<1-2>
\newcommand{\updownarrowlarge}{\multirow{3}{*}{\scalebox{1}[3]{$\updownarrow$}}}
	Also wurde das mittlere B vom unteren b empfangen?\mybreak

\only<1-2>{\begin{columns}
	\begin{column}{4cm}
		\begin{tabular}{cc|c|ccccc}
			& A & b & c & A\\
			& \only<2>{\updownarrowlarge} & $\updownarrow$\\
			& & B & a & B & c & B\\
			& & $\updownarrow$\\  
			C & a & b & C
		\end{tabular}
		
	\end{column}

	\begin{column}{5cm}
		\begin{tabular}{cc|c|cccccc}
			& A &   & b & c & A &\\
			& \only<2>{\updownarrowlarge} & \\
			&   & B & a & B & c & B\\
			& & $\updownarrow$\\  
			C & a & b & C
		\end{tabular}
	\end{column}
\end{columns}\mybreak}


\only<3>{\begin{columns}
	\begin{column}{4cm}
		\begin{tabular}{cc|c|ccccc}
			& A & b & c & A\\
			& \updownarrowlarge & $\updownarrow$ & \updownarrowlarge & $\updownarrow$\\
			&   & B &   & a & B & c & B\\
			& & $\updownarrow$\\  
			C & a & b & C
		\end{tabular}
	\end{column}

	\begin{column}{5cm}
		\begin{tabular}{cc|c|cccccc}
			& A &   & b & c & A &\\
			& \updownarrowlarge &  && \updownarrowlarge & $\updownarrow$\\
			&   & B &   &   & a & B & c & B\\
			& & $\updownarrow$\\  
			C & a & b &   & C
		\end{tabular}
	\end{column}
\end{columns}\mybreak}

\only<3>{	Nein!}
	
\end{frame}

\section{Alternativer Algorithmus}


\begin{frame}
Gibt es eine Lösung zu einer Teilmenge der Logs?\mybreak

Betrachte Schnitte:\mybreak
\only<1>{
	\[
	\begin{array}{cccccccc}
	\ldots & e_1 & \ldots\\
	\ldots & e_2 & \ldots\\[-3mm]
	       &     &  \underbrace{}_{\text{lösbar?}}
	\end{array}
	\]
}
\only<2>{
	\[
	\begin{array}{cccccccc}
	\ldots & e_1 & S & \ldots & S & e'_1 & \ldots\\
	\ldots & e_2 & S & \ldots & S & e'_2 & \ldots\\
	\end{array}
	\]
}
\only<3>{	
	\[
	\begin{array}{cccccccc}
	\ldots & e_1 &  \ldots  \\
	\ldots & e_2 &  \ldots\\
	\ldots & e_3 &  \ldots\\
	\ldots & e_4 &  \ldots  \\[-3mm]
&     &  \underbrace{}_{\text{lösbar?}}	
	\end{array}
	\]
}
\only<4-5>{	
	\[
	\begin{array}{cccccccc}
	\ldots & e_1 &  \ldots  \\
	\ldots & e_2 & S & \ldots & S & e'_2 & \ldots\\
	\ldots & e_3 & S & \ldots & S & e'_3 & \ldots\\
	\ldots & e_4 &  \ldots  \\
	\end{array}
	\]
}

\only<5>{
	Dasselbe für links $\Rightarrow$ Algorithmus für uSRC in $O(l^2\cdot n^{l})$}
\end{frame}

	
%\newcount\opaqueness
%\newdimen\offset
%\begin{frame}
%	\frametitle{Flying Theorems (You Really Shouldn’t!)}
%	\animate<2-14>
%	\animatevalue<1-15>{\opaqueness}{100}{0}
%	\animatevalue<1-15>{\offset}{0cm}{-5cm}
%	\begin{colormixin}{\the\opaqueness!averagebackgroundcolor}
%		\hskip\offset
%		\begin{minipage}{\textwidth}
%			\begin{theorem}
%				This theorem flies out.
%			\end{theorem}
%		\end{minipage}
%	\end{colormixin}
%	\animatevalue<1-15>{\opaqueness}{0}{100}
%	\animatevalue<1-15>{\offset}{-5cm}{0cm}
%	\begin{colormixin}{\the\opaqueness!averagebackgroundcolor}
%		\hskip\offset
%		\begin{minipage}{\textwidth}
%			\begin{theorem}
%				This theorem flies in.
%			\end{theorem}
%		\end{minipage}
%	\end{colormixin}
%\end{frame}
%
\begin{frame}
\begin{center}
	\animate<1-3>
	Danke für Ihre Aufmerksamkeit.\mybreak
	Fragen?
\end{center}	
\end{frame}

\begin{frame}[fragile]
	\animate<1-16>
	\setlength{\unitlength}{0.9cm}
	\newcommand{\tabtemp}{\phantom{x}}
	\newcommand{\tabuarrow}[1][\tabtemp]{#1} %\begin{picture}(0,0)\put(0.15,0.25){\vector(0,1){0.5}}\end{picture}#1}
	\newcommand{\tablarrow}[1][\tabtemp]{\begin{picture}(0,0)\put(0,0.15){\vector(-1,0){0.5}}\end{picture}#1}
	\newcommand{\tabluarrow}[1][\tabtemp]{\tablarrow[]\tabuarrow[#1]}
	\newcommand{\tabnwarrow}[1][\tabtemp]{\begin{picture}(0,0)\put(0,0.25){\vector(-1,1){0.5}}\end{picture}#1}
	\newcommand{\tabdarrow}[1][\tabtemp]{\begin{picture}(0,0)\put(0.15,0){\vector(0,-1){0.5}}\end{picture}#1}
	\newcommand{\tabrarrow}[1][\tabtemp]{\begin{picture}(0,0)\put(0.25,0.1){\vector(1,0){0.5}}\end{picture}#1}
	\newcommand{\tabsearrow}[1][\tabtemp]{\begin{picture}(0,0)\put(0.25,0.1){\vector(1,-1){0.5}}\end{picture}#1}
	\newcommand{\tabrdarrow}[1][\tabtemp]{\tabrarrow[]\tabdarrow[#1]}
	\only<1-9>{
		Anhang - Dynamische Programmierung für den Spezialfall zweier Logs:\mybreak
		
		Eingabe:\mybreak	
		\begin{center}
			\begin{tabular}{cccccc}
				B & c & a & B &   & \\
				C & A & C & b & C & A
			\end{tabular} \mybreak
		\end{center}
		
		Geschrieben als Matrix:\mybreak}
		\only<10-16>{Rechtslösbar:\mybreak}
		\begin{center}
			\only<1-2>{
				\begin{tabular}{c||cccc|cc}
					&C	&A 	& C	& b	& C	& A \\ \hline
					\hline B & 	&  	&  	& ?	&  	&  \\ 
					c &?	&  	&?	&  	& ?	&  \\ 
					\hline a & 	&? 	&  	&  	&  	& \alt<1>{?}{\alert{$*$}} \\ 
					\hline B & 	&  	&  	& \alt<1>{?}{\alert{$*$}}	&  	&  \\ 
					
				\end{tabular} 
			}
			\only<3->{
				%\newsavebox{\tabluarrowbox}
				%\savebox{\tabluarrowbox}{\vspace*{-0.5cm}\hspace*{-0.5cm}$\nwarrow$}
				%\newcommand{\tabluarrow}{\usebox{\tabluarrowbox}}
				%\newsavebox{\luarrowbox}
				%\savebox{\luarrowbox}{\vspace*{-0.5cm}\hspace*{-0.5cm}$\nwarrow$}
				%\newcommand{\tablarrow}{\hspace*{-0.5cm}$\longleftarrow$}
				%\newcommand{\textlonguparrow}{\scalebox{1}[2]{$\uparrow$}}
				%\newsavebox{\tabuarrowbox}
				%\savebox{\tabuarrowbox}{\vspace*{1cm}$\uparrow$}
				\begin{tabular}{c||cccc|cc}
					&C	&A 	& C	& b	& C	& A \\ \hline
					\hline B & 	&\tablarrow  	& \tablarrow 	&\tablarrow[\temporal<5>{?}{\alert{$+$}}{$+$}]	&  	&  \\ 
					c &\tabuarrow[\temporal<8>{?}{\alert{$+$}}{$+$}]	&  \tabluarrow	&\tabluarrow[\temporal<6>{?}{\alert{$0$}}{$0$}]	&  	&\tabuarrow[]\tabnwarrow[\temporal<4>{?}{\alert{$+$}}{$+$}]	&  \\ 
					\hline a & 	&\tabnwarrow[]\tablarrow[\temporal<7>{?}{\alert{$+$}}{$+$}] 	&  	& 	\tabnwarrow&  &\tablarrow[]\tabnwarrow[$*$] \\ 
					\hline B & 	&\tablarrow   &\tabnwarrow[]\tablarrow  &\tablarrow[$*$]	&  	&  \\ 
				\end{tabular} 
			}
		\end{center}	
		
		\only<10-16>{Linkslösbar:\mybreak
			\begin{center}
				
				
				
				\begin{tabular}{c||ccc|ccc}
					&C	&A 	& C	& b	& C	& A \\ \hline
					\hline B & 	&  	& 	&\tabsearrow[]\tabrarrow[\temporal<11>{?}{\alert{$*$}}{$*$}]	&  \tabrarrow	&  \\ 
					\hline  c &\tabsearrow[]\tabrarrow[\temporal<11>{?}{\alert{$*$}}{$*$}]	&  \tabrarrow	&\temporal<11>{?}{\alert{$*$}}{$*$}	&  	&\tabrarrow[]\tabsearrow[\temporal<14>{?}{\alert{$+$}}{$+$}]	&  \\ 
					\hline a & 	&\tabrdarrow[\temporal<12>{?}{\alert{$+$}}{$+$}] 	& \tabdarrow[] \tabsearrow 	& 	&  &\tabdarrow[\temporal<15>{?}{\alert{$+$}}{$+$}] \\ 
					B & 	& \tabrarrow  &&\tabrarrow[\temporal<13>{?}{\alert{$+$}}{$+$}]	&  \tabrarrow	&  \\ 
				\end{tabular} 
				
				%\hspace{0.5cm}
				%$\Longrightarrow$
				%\hspace{0.5cm}
				%\begin{minipage}{7cm}
				%	\begin{tabular}{c||ccc|ccc}
				%		&C	&A 	& C	& b	& C	& A \\ \hline
				%		\hline B & 	&  	& 	&\tabsearrow[]\tabrarrow[*]	&  \tabrarrow	&  \\ 
				%		\hline  c &\tabsearrow[]\tabrarrow[*]	&  \tabrarrow	&*	&  	&\tabrarrow[]\tabsearrow[+]	&  \\ 
				%		\hline a & 	&\tabrdarrow[+] 	& \tabdarrow[] \tabsearrow 	& 	&  &\tabdarrow[+] \\ 
				%		B & 	& \tabrarrow  &&\tabrarrow[+]	&  \tabrarrow	&  \\ 
				%	\end{tabular} 
				%\end{minipage}
				%	
			\end{center}}
		\end{frame}
		
\end{document}
